% GENERATED FILE. Edit canonical lessons, then run npm run generate:books.
% canonical-source-hash: fnv1a64:4ccb76db92c1c387
% canonical-lessons: JA-C466-shodou, JA-C466-ikebana, JA-C466-sadou, JA-C466-origami, JA-C466-geijutsu

\chapter{The Traditional Arts}
\label{ch:ja-the-traditional-arts}
\hlchaptermodality{466}

\begin{chapteropening}
\textbf{By the end of this chapter:} \emph{I can name some traditional Japanese arts.}

Complete the last lesson of chapter 466: I can name some traditional Japanese arts.
\end{chapteropening}

\section[\texorpdfstring{\ja{しょどう} (shodō)}{しょどう (shodō)}]{\ja{しょどう} (shodō) — calligraphy}

\label{lesson:JA-C466-shodou}



\begin{quote}
Before the new one: say the Japanese for a painter, then the Japanese for fine art.
\end{quote}

\subsection*{You'll want to know: \ja{しょどう}}
\textbf{\ja{しょどう}} — \emph{shodō} — \textquotedblleft{}calligraphy\textquotedblright{}.

\begin{quote}
\textbf{\ja{しゅうまつに} \ja{しょどうを} \ja{ならっています。}} — \emph{shūmatsu ni shodō o naratte imasu} — \textquotedblleft{}I take calligraphy lessons at the weekend.\textquotedblright{}
\end{quote}

\begin{cousinweb}
\textbf{\ja{しょどう}} (\emph{shodō}) is Sino-Japanese (kango): two Chinese characters read with their on'yomi, readings borrowed from Chinese. They are \emph{sho}, \textquotedblleft{}write\textquotedblright{} (Middle Chinese \emph{syo}) and \emph{dō}, \textquotedblleft{}way\textquotedblright{} (\emph{dawX}) (KANJIDIC2; Middle Chinese in Baxter's transcription).
\end{cousinweb}

\begin{grammarlens}[title={What you've built}]
One more word for describing what you lived and felt.
\end{grammarlens}

\subsection*{Your turn}
\begin{itemize}
\raggedright
  \item \emph{Say it:} \emph{shodō}
  \item \emph{Say it:} \emph{shodō}, once more
  \item \emph{Say it:} \emph{shūmatsu ni shodō o naratte imasu} — I take calligraphy lessons at the weekend.
  \item \emph{Recall:} say the Japanese for a painter, then the Japanese for fine art, then say \emph{shodō} again
\end{itemize}

\begin{tcolorbox}[breakable,colback=teal!4,colframe=teal!35!black,title={Before you move on}]
What does \ja{しょどう} mean? (\textquotedblleft{}Calligraphy\textquotedblright{}.) Say it once more.
\end{tcolorbox}
\section[\texorpdfstring{\ja{いけばな} (ikebana)}{いけばな (ikebana)}]{\ja{いけばな} (ikebana) — flower arranging}

\label{lesson:JA-C466-ikebana}



\begin{quote}
Before the new one: say the Japanese for fine art, then the Japanese for calligraphy.
\end{quote}

\subsection*{You'll want to know: \ja{いけばな}}
\textbf{\ja{いけばな}} — \emph{ikebana} — \textquotedblleft{}flower arranging\textquotedblright{}.

\begin{quote}
\textbf{\ja{ははは} \ja{いけばなが} \ja{じょうずです。}} — \emph{haha wa ikebana ga jōzu desu} — \textquotedblleft{}My mother is good at flower arranging.\textquotedblright{}
\end{quote}

\begin{cousinweb}
\textbf{\ja{いけばな}} (\emph{ikebana}) is a native Japanese compound (wago), inherited rather than borrowed from Chinese: \emph{ike}, \textquotedblleft{}keeping alive (ikeru)\textquotedblright{} and \emph{bana}, \textquotedblleft{}flower (hana)\textquotedblright{} (JMdict, KANJIDIC2).
\end{cousinweb}

\begin{grammarlens}[title={What you've built}]
One more word for describing what you lived and felt.
\end{grammarlens}

\subsection*{Your turn}
\begin{itemize}
\raggedright
  \item \emph{Say it:} \emph{ikebana}
  \item \emph{Say it:} \emph{ikebana}, once more
  \item \emph{Say it:} \emph{haha wa ikebana ga jōzu desu} — My mother is good at flower arranging.
  \item \emph{Recall:} say the Japanese for fine art, then the Japanese for calligraphy, then say \emph{ikebana} again
\end{itemize}

\begin{tcolorbox}[breakable,colback=teal!4,colframe=teal!35!black,title={Before you move on}]
What does \ja{いけばな} mean? (\textquotedblleft{}Flower arranging\textquotedblright{}.) Say it once more.
\end{tcolorbox}
\section[\texorpdfstring{\ja{さどう} (sadō)}{さどう (sadō)}]{\ja{さどう} (sadō) — the tea ceremony}

\label{lesson:JA-C466-sadou}



\begin{quote}
Before the new one: say the Japanese for calligraphy, then the Japanese for flower arranging.
\end{quote}

\subsection*{You'll want to know: \ja{さどう}}
\textbf{\ja{さどう}} — \emph{sadō} — \textquotedblleft{}the tea ceremony\textquotedblright{}.

\begin{quote}
\textbf{\ja{さどうの} \ja{せんせいに} \ja{あいました。}} — \emph{sadō no sensei ni aimashita} — \textquotedblleft{}I met a teacher of the tea ceremony.\textquotedblright{}
\end{quote}

\begin{cousinweb}
\textbf{\ja{さどう}} (\emph{sadō}) is Sino-Japanese (kango): two Chinese characters read with their on'yomi, readings borrowed from Chinese. They are \emph{sa}, \textquotedblleft{}tea\textquotedblright{} (Middle Chinese \emph{drae}) and \emph{dō}, \textquotedblleft{}way\textquotedblright{} (\emph{dawX}) (KANJIDIC2; Middle Chinese in Baxter's transcription).
\end{cousinweb}

\begin{grammarlens}[title={What you've built}]
One more word for describing what you lived and felt.
\end{grammarlens}

\subsection*{Your turn}
\begin{itemize}
\raggedright
  \item \emph{Say it:} \emph{sadō}
  \item \emph{Say it:} \emph{sadō}, once more
  \item \emph{Say it:} \emph{sadō no sensei ni aimashita} — I met a teacher of the tea ceremony.
  \item \emph{Recall:} say the Japanese for calligraphy, then the Japanese for flower arranging, then say \emph{sadō} again
\end{itemize}

\begin{tcolorbox}[breakable,colback=teal!4,colframe=teal!35!black,title={Before you move on}]
What does \ja{さどう} mean? (\textquotedblleft{}The tea ceremony\textquotedblright{}.) Say it once more.
\end{tcolorbox}
\section[\texorpdfstring{\ja{おりがみ} (origami)}{おりがみ (origami)}]{\ja{おりがみ} (origami) — paper folding; origami}

\label{lesson:JA-C466-origami}



\begin{quote}
Before the new one: say the Japanese for flower arranging, then the Japanese for the tea ceremony.
\end{quote}

\subsection*{You'll want to know: \ja{おりがみ}}
\textbf{\ja{おりがみ}} — \emph{origami} — \textquotedblleft{}paper folding; origami\textquotedblright{}.

\begin{quote}
\textbf{\ja{むすめは} \ja{おりがみが} \ja{じょうずです。}} — \emph{musume wa origami ga jōzu desu} — \textquotedblleft{}My daughter is good at paper folding.\textquotedblright{}
\end{quote}

\begin{cousinweb}
\textbf{\ja{おりがみ}} (\emph{origami}) is a native Japanese compound (wago), inherited rather than borrowed from Chinese: \emph{ori}, \textquotedblleft{}folding (oru)\textquotedblright{} and \emph{gami}, \textquotedblleft{}paper (kami)\textquotedblright{} (JMdict, KANJIDIC2).
\end{cousinweb}

\begin{grammarlens}[title={What you've built}]
One more word for describing what you lived and felt.
\end{grammarlens}

\subsection*{Your turn}
\begin{itemize}
\raggedright
  \item \emph{Say it:} \emph{origami}
  \item \emph{Say it:} \emph{origami}, once more
  \item \emph{Say it:} \emph{musume wa origami ga jōzu desu} — My daughter is good at paper folding.
  \item \emph{Recall:} say the Japanese for flower arranging, then the Japanese for the tea ceremony, then say \emph{origami} again
\end{itemize}

\begin{tcolorbox}[breakable,colback=teal!4,colframe=teal!35!black,title={Before you move on}]
What does \ja{おりがみ} mean? (\textquotedblleft{}Paper folding; origami\textquotedblright{}.) Say it once more.
\end{tcolorbox}
\section[\texorpdfstring{\ja{げいじゅつ} (geijutsu)}{げいじゅつ (geijutsu)}]{\ja{げいじゅつ} (geijutsu) — the arts; art}

\label{lesson:JA-C466-geijutsu}



\begin{quote}
Before the new one: say the Japanese for the tea ceremony, then the Japanese for paper folding.
\end{quote}

\subsection*{You'll want to know: \ja{げいじゅつ}}
\textbf{\ja{げいじゅつ}} — \emph{geijutsu} — \textquotedblleft{}the arts; art\textquotedblright{}.

\begin{quote}
\textbf{\ja{まちの} \ja{げいじゅつの} \ja{まつりに} \ja{いきました。}} — \emph{machi no geijutsu no matsuri ni ikimashita} — \textquotedblleft{}I went to the town's arts festival.\textquotedblright{}
\end{quote}

\begin{cousinweb}
\textbf{\ja{げいじゅつ}} (\emph{geijutsu}) is Sino-Japanese (kango): two Chinese characters read with their on'yomi, readings borrowed from Chinese. They are \emph{gei}, \textquotedblleft{}art\textquotedblright{} (Middle Chinese \emph{ngjiejH}) and \emph{jutsu}, \textquotedblleft{}art\textquotedblright{} (\emph{zywit}) (KANJIDIC2; Middle Chinese in Baxter's transcription).
\end{cousinweb}

\begin{grammarlens}[title={What you've built}]
One more word for describing what you lived and felt.
\end{grammarlens}

\subsection*{Your turn}
\begin{itemize}
\raggedright
  \item \emph{Say it:} \emph{geijutsu}
  \item \emph{Say it:} \emph{geijutsu}, once more
  \item \emph{Say it:} \emph{machi no geijutsu no matsuri ni ikimashita} — I went to the town's arts festival.
  \item \emph{Recall:} say the Japanese for the tea ceremony, then the Japanese for paper folding, then say \emph{geijutsu} again
\end{itemize}

\begin{tcolorbox}[breakable,colback=teal!4,colframe=teal!35!black,title={Before you move on}]
What does \ja{げいじゅつ} mean? (\textquotedblleft{}The arts; art\textquotedblright{}.) Say it once more.
\end{tcolorbox}
